\documentclass{article}
\usepackage{graphicx} % Required for inserting images
\usepackage{amsmath} 

\title{Cálculo Numérico - Aspectos Teóricos e Computacionais \\ \large Capítulo 1}
\author{Daniel Santos}

\begin{document}

\maketitle

\section*{Exercícios}
\begin{enumerate}
\item Converta os seguintes números decimais para sua forma binária:\\
$x = 37$ \hfil
$y = 2345$ \hfil
$z = 0.1217$

\item Converta os seguintes números binários para sua forma decimal:\\
$x = (101101)_2$  \hfill  $y = (110101011)_2$ \hfill \mbox{} \\
$z = (0.01101)_2$   \hfill  $w = (0.111111101)_2$ \hfill \mbox{}

\item Seja um sistema de aritmética de ponto flutuante de quatro\\ dígitos, base decimal e com acumulador de precisão dupla. Dados os números: \\
$x=0.7237 \times 10^{4}$\quad
$y = 0.2145 \times 10^{-3}$ \quad
e \quad $z=0.2585 \times 10^{1}$ \\
efetue as seguintes operações e obtenha o erro relativo no resultado,\\ supondo que x, y e z estão exatamente representados:
\begin{enumerate}
    \item $x+y+z$
    \item $x-y-z$
    \item $x/y$
    \item $(xy)/z$
    \item $x(y/z)$
\end{enumerate}

\end{enumerate}


\section*{Resolução}
\begin{enumerate}
\item 
\begin{enumerate}
    \item 
        \begin{align*}
        &37 \\
        &N_0 = 37 = 2\cdot18 + 1 \xrightarrow{} a_0 = 1 \\ 
        &N_1 = 18 = 2\cdot9 + 0 \xrightarrow{} a_1 = 0  \\
        &N_2 = 9 = 2\cdot4 + 1 \xrightarrow{} a_2 = 1  \\
        &N_3 = 4 = 2\cdot2 + 0 \xrightarrow{} a_3 = 0  \\
        &N_4 = 2 = 2\cdot1 + 0 \xrightarrow{} a_4 = 0  \\
        &N_5 = 1 = 2\cdot0 + 1 \xrightarrow{} a_5 = 1  \\
        &(37)_10 = (100101)_2
        \end{align*}
    \item 
        \begin{align*}
        &2345 \\
        &N_0 = 2345 = 2\cdot1172 + 1 \xrightarrow{} a_0 = 1 \\ 
        &N_1 = 1172 = 2\cdot586 + 0 \xrightarrow{} a_1 = 0 \\ 
        &N_2 = 586 = 2\cdot293 + 0  \xrightarrow{} a_2 = 0 \\ 
        &N_3 = 293 = 2\cdot146 + 1 \xrightarrow{} a_3 = 1 \\ 
        &N_4 = 146 = 2\cdot73 + 0  \xrightarrow{} a_4 = 0 \\ 
        &N_5 = 73 = 2\cdot36 + 1 \xrightarrow{} a_5 = 1 \\ 
        &N_6 = 36 = 2\cdot18 + 0 \xrightarrow{} a_6 = 0 \\ 
        &N_7 = 18 = 2\cdot9 + 0 \xrightarrow{} a_7 = 0 \\ 
        &N_8 = 9 = 2\cdot4 + 1 \xrightarrow{} a_8 = 1  \\
        &N_9 = 4 = 2\cdot2 + 0 \xrightarrow{} a_9 = 0  \\
        &N_10 = 2 = 2\cdot1 + 0 \xrightarrow{} a_10 = 0  \\
        &N_11 = 1 = 2\cdot0 + 1 \xrightarrow{} a_11 = 1  \\
        &(2345)_10 = (100100101001)_2
        \end{align*}
    \item 
        \begin{align*}
        &0.1217 \\
        & k=1 \quad  2r_1 = 0.2434 => d_1 = 0 \quad r_2 = 0.2434 \\ 
        & k=2 \quad  2r_2 = 0.4868 => d_2 = 0 \quad r_3 = 0.4868  \\
        & k=3 \quad  2r_3 = 0.9736 => d_3 = 0 \quad r_4 = 0.9736  \\
        & k=4 \quad  2r_4 = 1.9472 => d_4 = 1 \quad r_5 = 0.9472  \\
        & k=5 \quad  2r_5 = 1.8944 => d_5 = 1 \quad r_6 = 0.8944  \\
        & k=6 \quad  2r_6 = 1.7888 => d_6 = 1 \quad r_7 = 0.7888  \\
        & k=7 \quad  2r_7 = 1.5776 => d_7 = 1 \quad r_8 = 0.5776  \\
        & k=8 \quad  2r_8 = 1.1552 => d_8 = 1 \quad r_9 = 0.1552  \\
        & k=9 \quad  2r_9 = 0.3104 => d_9 = 0 \quad r_{10} = 0.3104  \\
        & k=10 \quad  2r_{10} = 0.6208 => d_{10} = 0 \quad r_{11} = 0.6208  \\
        & k=11 \quad  2r_{11} = 1.2416 => d_{11} = 1 \quad r_{12} = 0.2416  \\
        & k=12 \quad  2r_{12} = 0.4832 => d_{12} = 0 \quad r_{13} = 0.4832  \\
        & k=13 \quad  2r_{13} = 0.9664 => d_{13} = 0 \quad r_{14} = 0.9664  \\
        & k=14 \quad  2r_{14} = 1.9328 => d_{14} = 1 \quad r_{15} = 0.9328  \\
        & k=15 \quad  2r_{15} = 1.8656 => d_{15} = 1 \quad r_{16} = 0.8656  \\
        & k=16 \quad  2r_{16} = 1.7312 => d_{16} = 1 \quad r_{17} = 0.7312  \\
        & k=17 \quad  2r_{17} = 1.4624 => d_{17} = 1 \quad r_{18} = 0.4624  \\
        & k=18 \quad  2r_{18} = 0.9248 => d_{18} = 0 \quad r_{19} = 0.9248  \\
        & k=19 \quad  2r_{19} = 1.8496 => d_{19} = 1 \quad r_{20} = 0.8496  \\
        & k=20 \quad  2r_{20} = 1.6992 => d_{20} = 1 \quad r_{21} = 0.6992  \\
        & k=21 \quad  2r_{21} = 1.3984 => d_{21} = 1 \quad r_{22} = 0.3984  \\
        & k=22 \quad  2r_{22} = 0.7968 => d_{22} = 0 \quad r_{23} = 0.7968  \\
        & k=23 \quad  2r_{23} = 1.5936 => d_{23} = 1 \quad r_{24} = 0.5936  \\
        & k=24 \quad  2r_{24} = 1.1872 => d_{24} = 1 \quad r_{25} = 0.1872  \\
        \end{align*}

        E assim segue infinitamente, não existindo representação exata.
\end{enumerate}

\item 
\begin{enumerate}
    \item 
        \begin{align*}
        &(101101)_2 \\
        &= 1\cdot2^5 + 0\cdot2^4 + 1\cdot2^3 + 1\cdot2^2 + 0\cdot2^1 + 1\cdot2^0 \\
        &= 32 + 0 + 8 + 4 + 0 + 1 = 45
        \end{align*}
    \item 
        \begin{align*}
        &(110101011)_2 \\
        &= 1\cdot2^8 +1\cdot2^7 +0\cdot2^6 +1\cdot2^5 + 0\cdot2^4 + 1\cdot2^3 \\
        &+ 0\cdot2^2 + 1\cdot2^1 + 1\cdot2^0 \\
        &= 256 + 128 + 0 + 32 + 0 + 8 + 0 +2 + 1 \\
        &= 427
        \end{align*}
    \item 
        \begin{align*}
        &(0.01101)_2 \\
        &=0\cdot2^{-1} + 1\cdot2^{-2} + 1\cdot2^{-3} + 0\cdot2^{-4} + 1\cdot2^{-5} \\
        &=0 + 0.25 + 0.125 + 0 + 0.03125 \\
        &=0.40625
        \end{align*}
    \item 
        \begin{align*}
        &(0.111111101)_2 \\
        &=1\cdot2^{-1} + 1\cdot2^{-2} + 1\cdot2^{-3} + 1\cdot2^{-4}  \\
        &+ 1\cdot2^{-5} + 1\cdot2^{-6} + 1\cdot2^{-7} + 0\cdot2^{-8} + 1\cdot2^{-9} \\
        &= 0.5 + 0.25 + 0.125 + 0.0625 + 0.03125 + 0.015625 \\
        &+ 0.0078125 + 0 + 0.001953125 \\
        &=0.994140625 
        \end{align*}
\end{enumerate}

\item
\begin{enumerate}
    \item 
    Valor de y desconsiderado por conta do acumulador de precisão \\dupla ter precisão apenas até a 8ª cada decimal.
    \begin{align*}
        &\overline{x + y} \\
        &=0.7237 \times 10^{4} + 0.2145 \times 10^{-3}\\
        &=0.7237 \times 10^{4} + 0.000000002145 \times 10^{4}\\
        &=0.7237 \times 10^{4} \\
        &\overline{(x+y)+z} \\
        &0.7237 \times 10^{4} + 0.2585 \times 10^{1}\\
        &=0.7237 \times 10^{4} + 0.0002585 \times 10^{4}\\
        &=0.7239585 \times 10^{4}
    \end{align*}
    
    $\overline{x+y+z} = 0.7239 \times 10^{4}$, no truncamento. \\
    $\overline{x+y+z} = 0.7240 \times 10^{4}$, no arredondamento. \\

    O valor real seria $0.72395850214\times 10^{4}$. Vamos calcular o erro relativo.\\Para simplificar, considerarei apenas o arredondamento.
    \begin{align*}
        EA_{x+y+z} &= |(x+y+z) - \overline{x+y+z}|\\
        &=|0.72395850214\times 10^{4} - 0.7240 \times 10^{4}|\\
        &=0.00004149786 \times10^{4} = 0.4149786\\
        ER_{x+y+z} &= \frac{EA_{x+y+z}}{\overline{x+y+z}}\\
        &=\frac{0.4149786}{0.7240 \times 10^{4}}\\
        &\approx 0.00005731748
    \end{align*}
    \item 
    Valor de y desconsiderado por conta do acumulador de precisão \\dupla ter precisão apenas até a 8ª cada decimal.
    \begin{align*}
        &\overline{x - y} \\
        &=0.7237 \times 10^{4} - 0.2145 \times 10^{-3}\\
        &=0.7237 \times 10^{4} - 0.000000002145 \times 10^{4}\\
        &=0.7237 \times 10^{4} \\
        &\overline{(x-y)-z} \\
        &0.7237 \times 10^{4} - 0.2585 \times 10^{1}\\
        &=0.7237 \times 10^{4} - 0.0002585 \times 10^{4}\\
        &=0.7234415 \times 10^{4}
    \end{align*}

    $\overline{x-y-z} = 0.7234 \times 10^{4}$, no arredondamento e truncamento. \\

    O valor real seria $7237-0.0002145-2.585=7234.4147855$. Vamos calcular o erro relativo.
    \begin{align*}
        EA_{x-y-z} &= |(x-y-z) - \overline{x-y-z}|\\
        &=|7234.4147855 - 7234|\\
        &=0.4147855\\
        ER_{x+y+z} &= \frac{EA_{x+y+z}}{\overline{x+y+z}}\\
        &=\frac{0.4147855}{7234.4147855}\\
        &\approx 5.73\times 10^{-5}
    \end{align*}
    \item 
    \begin{align*}
        &\overline{x / y} \\
        &=(0.7237 \times 10^{4}) / (0.2145 \times 10^{-3})\\
        &=(0.7237/0.2145) \times 10^{7} \\
        &=3.37389277389\times 10^{7}\\
        &=0.337389277389\times 10^{8}\\
    \end{align*}

    $\overline{x/y} = 0.3373\times 10^{8}$, no truncamento. \\
    $\overline{x/y} = 0.3374\times 10^{8}$, no arredondamento. \\
    
    O valor real seria $7237/0.0002145 \approx 33738927.7389$.\\

    \begin{align*}
        EA_{x/y} &= |(x/y) - \overline{x/y}|\\
        &=|33740000 - 33738927.7389|\\
        &=1072.2611\\
        ER_{x/y} &= \frac{EA_{x/y}}{\overline{x/y}}\\
        &=\frac{1072.2611}{33738927.7389}\\
        &\approx 3.17\times 10^{-5}
    \end{align*}
    \item 
    \begin{align*}
        &\overline{(xy)}\\
        &=0.7237 \times 10^{4} \cdot 0.2145 \times 10^{-3}\\
        &=(0.7237\cdot0.2145) \times 10^{1}\\
        &=0.15523365 \times 10^{1} \\
        &=0.1552 \times 10^{1} \\
        &\overline{(xy)/z}\\
        & (0.1552 \times 10^{1}) / (0.2585 \times 10^{1}) \\
        & 0.60038684719 \times 10^{0}
    \end{align*}
    
    $\overline{(xy)/z} = 0.6003\times 10^{0}$, no truncamento. \\
    $\overline{(xy)/z} = 0.6004\times 10^{0}$, no arredondamento. \\

    O valor real seria $\approx 0.60051$.\\

    \begin{align*}
        ER_{(xy)/z} &= \frac{EA_{(xy)/z}}{\overline{(xy)/z}}\\
        &=\frac{0.60051 - 0.6004}{0.60051}\\
        &\approx 1.83\times 10^{-4}
    \end{align*}    
    \item 
        \begin{align*}
        &\overline{(y/z)}\\
        &=0.2145 \times 10^{-3}/0.2585 \times 10^{1}\\
        &= 0.82978723404 \times 10^{-2}\\
        &=0.8297 \times 10^{-2} \\
        &\overline{(x(y/z)}\\
        & 0.8297 \times 10^{-2} \cdot 0.7237 \times 10^{4}\\
        & 0.60045389 \times 10^{2}
    \end{align*}
    
    $\overline{(x(yz)} = 0.6004\times 10^{2}$, no truncamento. \\
    $\overline{((x(yz)} = 0.6005\times 10^{2}$, no arredondamento. \\

    \begin{align*}
        ER_{(x(y/z)} &= \frac{EA_{(x(y/z)}}{\overline{(x(y/z)}}\\
        &=\frac{|0.60051 - 0.6005|}{0.60051}\\
        &\approx 1.66 \times 10^{-5}
    \end{align*}    
\end{enumerate}
\end{enumerate}


\end{document}
